% Exact one-dimensional gauge-freedom example.
\clearpage
\section*{Inverse-problem laboratory: a zero residual is not a unique answer}
\noindent Suppose a measurement can determine only the difference
\(d=a-b\) between two unknown positions. The residual
\[
r(a,b)=a-b-d
\]
has Jacobian \(J=\begin{pmatrix}1&-1\end{pmatrix}\)
and Gauss--Newton normal matrix
\[
J^TJ=\begin{pmatrix}1&-1\\-1&1\end{pmatrix}.
\]
This matrix is positive semidefinite but singular:
\(J(1,1)^T=0\). Adding any common translation \(c\) to
both \(a\) and \(b\) leaves the measurement and residual unchanged.
The direction \((1,1)\) is a \emph{gauge freedom}, not noise.

Choosing \(b=0\) gives the representative \(a=d\).
That choice fixes coordinates; it does not create a new observation
or make the absolute location physically known.
An arbitrary diagonal regularizer might make the matrix
invertible, but it selects a preference rather than identifying
a previously unobservable quantity.

The same issue appears in pose estimation and reconstruction.
A collection of relative observations can determine scene geometry
up to a global rigid motion or similarity transform, depending on
the actual measurement model. Before interpreting a fitted pose,
count independent constraints, inspect Jacobian rank and distinguish
a true physical anchor from an imposed numerical gauge.

\paragraph{Working exercise.}
Derive the \(2\times2\) singular-value structure of the example.
Create three equally valid least-squares solutions for a fixed \(d\),
and plot their identical residual. Then fix one coordinate, recompute
the reduced problem and explain which information is gained versus
which ambiguity is merely named away. Extend the example to 2D
points with pairwise differences.
